\subsection{The complete generic pinned local input}
\label{sec:flag-generic-pinned-input}

Retain the concrete setting of the equal-word generic branch,
with the standing hypotheses and the actual \(\mathrm{Ctx}\),
\(\mathrm{MassDataTie}\), \(\mathrm{MassDataAdmissible}\) and
\(\mathrm{Pinned}\) data. The following supplies
\(\mathrm{StepLocGenPinnedG}\) in this setting, with no estimate
for the real orbital energy used in its proof.

\begin{lemma}[Exact scale and the actual normalized class square]
\label{lem:flag-pin-scale}
Let \(\nu_B\) be any of the original pinned real quotient
measure assignments. For a representable rational point \(w\)
let \(b(w)=\#\mathrm{FreeE}(w)\), the number of pairs \(i<j\)
belonging to one block of the point. There is a positive scale
\(s_{\rm pin}(\gamma)\le1\) on the real group, equal to one on
its regular locus, such that
\[
  s_{\rm pin}(\mathrm{gamW}(w))=g^{-fb(w)} .              \tag{GP1}
\]
Let \(\nu_L=s_{\rm pin}\nu_B\) denote multiplication of each
quotient measure by the scalar assigned to its point.
If the point is occupied by a nonzero generic term, then
\[
 \begin{split}
 J_2^L(w)&=g^{-2fb(w)}J_2^B(w)>0,\\
 \frac{\|\mathrm{grpGC}(c,w,\nu)\|^2}{J_2^L(w)}
       &=\mathrm{wt}(c)^2
             \prod_{v\in P_*}|A_v^0(d_v)|^2
 \end{split}
 \tag{GP2}
\]
for every nonempty generic key \((c,\nu)\) at this point.
The symbols \(J_2^L,J_2^B\) are the actual \(\mathrm{J2of}\)
with the two indicated measures; \(P_*\) is the finite
bad-place set of (GC3).
\end{lemma}
\begin{proof}
For any real group element that exactly represents a rational
point, the pattern of vanishing leading minors of its signed
cross matrix determines all the block ends: at a block
antidiagonal matrix with nonzero block scalars the nonzero
leading minors occur precisely at the block ends. Hence the
integer \(b(w)=\sum_j \binom{t_j}{2}\), where \(t_j\) are
the block sizes, is independent of the point representative.
In fact the full point is determined, since the real embeddings
of \(E\) and the fixed coordinate transport are injective.
If the real group element is regular, all these leading minors
are nonzero, and \(b(w)=0\). Set \(s_{\rm pin}(\gamma)\)
equal to \(g^{-fb(w)}\) on the exactly representable locus,
and one elsewhere. This is a well-defined scalar assignment;
no integration over the parameter of this assignment is needed.
Since \(g\ge1\), the asserted bounds hold globally, including
the real group elements not occurring in the arithmetic sum.
The pinned quotient conditions remain valid as quotient
conditions after rescaling: rescale the stabilizer Haar measure
in the opposite direction in the Weil formula. On the regular
locus neither measure changes.

At an occupied nonzero-Gram point its rational adapted frames
have the unique signed block-antidiagonal cross matrix described
in Section~\ref{sec:flag-nondegenerate-normal-count}. At each
infinite place they can be realized with the first flag fixed:
extend the isometry of the first adapted frame to a Witt
basis, and extend the second frame isometry in the same
nondegenerate ambient space. This gives a representative in
\(\mathrm{IsPtRep}\), with the original signed cross matrix
at every place. Thus \(\mathrm{J2of}\) uses its represented
branch, not the fallback zero. The orbital integral \(J_w\)
in (GC10) is precisely
\(\mathrm{scl}\,\mathrm{vol}(K_2)\) times the original
\(\mathrm{OO}\) with the pinned quotient measure. Linearity
of that integral in the measure gives the first equality
in (GP2). Its positivity is not inferred from positivity
of a sum of local phases: the individual term equation
preceding (GC10) shows that a nonzero term vector forces
\(J_w(\nu_B)\ne0\). Its finite phase has modulus one and
its arithmetic covolume and row weight are positive.
Now (GC10), which retains the complete sum over integral
returns in a rational key, gives the second equality in
(GP2) by dividing by the positive \(J_2^L(w)\).
For an empty key the original sum is zero and its
normalized square is zero.
\end{proof}

\begin{theorem}[Generic local supplier]
\label{thm:flag-generic-pinned-input}
Under the actual hypotheses above, one can take the uniform
constants
\[
                   C_1=4,\qquad C_0=1,\qquad C_{\rm St}=0,
                   \qquad e_1=2.                      \tag{GP3}
\]
In particular the weaker original allowances \(C_1=39,C_0=9\)
also suffice. With
\[
 {\cal A}=AA\,\rho^{2r(r+1)f(k-1)},\qquad
 N_{\rm gen}(w)=2{\cal A}\,
             H(w)^{10s^2+12s+1}\,\mathrm{nPt}(w),        \tag{GP4}
\]
the original data \(\mathrm{NOfG}(Y,u,\nu_L,\nu_B,N_{\rm gen})\)
satisfy \(\mathrm{LocalInputG}\) and \(\mathrm{GenMassInput}\)
on \(\mathrm{genT}\). The measure bound is in fact
\(J_2^L\le J_2^B\). The function (GP4) is defined by the
displayed expression on the original point space, even
at points not occupied by nonzero terms.
\end{theorem}
\begin{proof}
We describe the finite local witness required by the quantifiers
in \(\mathrm{LocalInputG}\). In particular, we do not assume
an infinite set of local data is a finite type.

Fix an occupied point, and fix an arbitrary finite test set
\(\mathcal F\) of original class and rational normal pairs.
Choose one nonzero generic term at the point, which exists
by \(\mathrm{IsPt}\) for \(\mathrm{genT}\).
Take its two rational adapted flags as in (GC1)--(GC10):
put their cross matrix in the prescribed signed point
normal form, and transport every other term at this point
to the same two rational flags. The transport exists by
adapted upper triangular changes followed by Hermitian
Witt extension on the nondegenerate joint span.
Changes on the joint span which preserve its two adapted
flags belong to the common unipotent, and changes on the
orthogonal complement belong to its definite unitary group.
Thus the two inner lattices and their level graph of a
nonempty key are well defined up to the complement action.
Here the class \(c=\mathrm{clsG}(t)\) records the row and its
scalar/corank-one equivalence; it therefore fixes the two
row lattices. The normal \(\nu=\mathrm{nmlOf}(t)\) is
exactly the rational double coset, not a selected
representative of it. Equality of the row and normal
is precisely the fixed-key comparison in (GC5);
refining by \(\mathrm{clsG}\) does not split this
key: under rational unipotent transport with the same
row the quotient maps and graph are unchanged, and
the scalar equality holds for the resulting common
nonzero-Gram flag data. This is the scalar/normal
comparison proved in the nondegenerate normal-count
section.

Use the local scaled flags and the reference datum of the
chosen term at each place. A place is bad if and only if
\(\mathrm{Hp}>1\): in (GL2) all valuations
\(v(\theta),a,d\) are nonnegative, and this height is a
power of the residue size. When it is one the place is
ordinary and all the block-end determinants and the
joint Gram determinant are units, exactly the
\(\mathrm{PrimeTo}\) condition. There are finitely many
bad places \(P_*\). Their height product is the prescribed
\(H(w)\), by (GC6). At every other place the integral
frame class is unique by (GL4), the joint lattice is
self-dual, and the lattice splits as the joint lattice
plus a self-dual complement. Integral self-dual lattices
of the given local complement are isometric: this is
the unramified Hermitian elementary-divisor classification
(or the lattice and its dual in the split case). Therefore
all occupied data outside \(P_*\) have the reference
complement orbit, and their single generic amplitude
has weight and phase one by (GC3). This last assertion
is about the term over a given datum; it does not assert
that the global definite class number is one.

At each \(v\in P_*\), let \(\mathcal D_v\) consist of the
finitely many complement-orbit representatives arising
from the nonempty tested keys, together with the chosen
reference datum. Choose the local weights and amplitudes
exactly as in (GL6)--(GL7):
\[
 \mathrm{wt}_v(d)=1/\operatorname{vol}(H(F_v)_d),\qquad
 a_Z(v,d)=A_v^0(d),\qquad
 nn_v=\sum_{d\ {\rm all\ local\ orbits}}\!
                \mathrm{wt}_v(d)m_d .
\]
The complement measures give the reference stabilizers
volume one. Take \(ngen_v\) to be the mass
\(ngen_v(C_v)\) in the local pairing class supplied
by the chosen generic term. This mass sums over all
local data; it is not truncated to \(\mathcal D_v\).
The sums are finite and positive as required by (GL7).
Choose \(V_0\) as in (GC4) and \(rw(c)\) as the squared
row weight (zero on an unused class). Assign to each
nonempty tested key its tuple of \(\mathcal D_v\) data.
Assign none to every empty or untested key. This defines
the optional datum map on the whole class/normal space.
For the unused \(a_N\) of the abstract record one can
use zero; \(\mathrm{LocalInputG}\) does not require
the nongeneric bound.

The positivity and residue-size conditions of
\(\mathrm{LocalNG.Basic}\) are now immediate from the
volume definitions and (GL2), since a bad height is
at least its residue size. The finite local sums satisfy
\(\mathrm{ZAt}\) by (GL7), because
\[
        e=s^2+3s\le4s^2\qquad(s\ge1).
\]
Likewise \(nn_v\le2H_v^{s^2+3s}ngen_v\)
implies the required
\(nn_v\le8H_v^{C_0s^2+4s}ngen_v\) with \(C_0=1\).
Every place in this record is bad; had a good place
been retained, its mass equality would hold by (GL7).
The proof of (GC4)--(GC5) applies to any finite subset
of keys with a fixed tuple of local data. Its result
is exactly \(\mathrm{Unfold}\):
\[
 \sum_{x\in \mathcal F'} rw(x.1)
          \le V_0\prod_{v\in P_*}\mathrm{wt}_v(d_v).
\]
The inequality is preserved by the optional-map
restriction to tested keys; no equality of sums over
all possible keys is required by this input.
Equation (GP2) proves the prescribed exact equality
of normalized class squares and
\(\mathrm{weightedGenericAmplitudeSq}\) on \(\mathcal F\),
including the keys assigned none.

Finally, the original mass of this finite witness is
\(M_{\mathcal F}=V_0\prod_{v\in P_*}ngen_v\).
The exact relative stabilizer formula (GC7), the
gluing mass calculation and Theorem
\ref{thm:flag-aa-normalization} give
\[
 0<M_{\mathcal F}\le
         2{\cal A}\,H(w)^{s(s-1)}\,\mathrm{nPt}(w).
\]
This is (AN16) for the pairing classes of the chosen
generic term. In particular the chosen reference
corrected span lattice is the actual one of that
term; it has the full required different and level
content, and its complement is definite. No
arbitrary choice of a local pairing class is made
without a global rational reference.

Since \(H(w)\ge1\) and \(s(s-1)\le10s^2+12s+1\),
we have \(M_{\mathcal F}\le N_{\rm gen}(w)\).
Enlarge only \(V_0\) by the positive multiplier
\(N_{\rm gen}(w)/M_{\mathcal F}\ge1\). This does not
change the local amplitudes, local masses, height,
or normalized class-square equalities. It preserves
\(\mathrm{Unfold}\) in its inequality direction and
makes \(\mathrm{ngenTot}=N_{\rm gen}(w)\) exactly,
as required by \(\mathrm{LocalInputG}\).
This harmless enlargement is not a change of the
Haar measures in the earlier class-factorization
or gluing arguments. It is only an enlargement of
the upper-bound record after those arguments.
The product over an empty \(P_*\) is interpreted as
one. Thus the construction works for every finite
test set, even an empty one.

By (GM10) the original \(\mathrm{nPt}(w)\) is
nonnegative (indeed positive on occupied points).
The defining exponent in \(\mathrm{GenMassInput}\)
is \(10s^2+6(2r+2)+1=10s^2+12s+1\).
Equation (GP4) therefore gives that input with
\(e_1=2\). The global scale satisfies
\(s_{\rm pin}^2\le1=\mathrm{b}^{0}\), proving the
claimed uniform \(C_{\rm St}=0\).
These constants are absolute and do not depend
on the dimension, field, point, or finite test set;
all standing thresholds may consequently be chosen
before the arithmetic datum. This completes the
generic pinned local input.
\end{proof}

The separate analytic point-sum input uses
\(\nu_B\), not \(\nu_L\). In particular this theorem
does not replace its stabilizer-root gauge by the
scaled gauge, and it does not make an estimate
for \(\mathrm{StepThmBPinnedG}\). The local theorem
also does not claim the nongeneric cancellation
saving \(q^{-(k-2r-5)}\). Those separate suppliers are
provided by the pointwise and exact-norm arguments below,
and by the second-type cancellation theorem, respectively.
